\bookchapter{Anhang C: Glossar der “Fachsprache”}

\begin{center}
\begin{tabularx}{\linewidth}{|X|X|X|}
\hline
\textbf{Der "große Speicher" des Smartphones} & \textbf{Hier liegen Fotos, Videos und Apps. 128G/256G beziehen sich darauf.} & \textbf{OLED / LCD} \tabularnewline
\hline
Zwei Bildschirmmaterialien & OLED farbenfroher und stromsparender; LCD augenfreundlicher und günstiger. & NFC \tabularnewline
Kabellose Nahfeldkommunikation & Macht aus deinem Smartphone eine Transitkarte, Zutrittskarte oder Bankkarte – die "Tappen"-Funktion. & OTA-Update \tabularnewline
Kabellose Systemaktualisierung & Das "System-Update", das du in den Einstellungen antippst – das ist es. & Root \tabularnewline
Android-Administratorrechte & Erlaubt das Ändern von Systemdateien, bringt aber auch Risiken. & BL-Sperre \tabularnewline
Bootloader-Sperre & Eine vom Hersteller gesetzte Sperre; erst nach dem Entsperren lässt sich flashen. & NPU \tabularnewline
Neuraler Prozessor & Ein auf KI-Aufgaben spezialisierter Chip, der die CPU um ein Vielfaches übertrifft und dabei weniger Strom verbraucht. & On-Device-KI \tabularnewline
Lokal auf dem Smartphone laufende KI & KI läuft direkt auf dem Smartphone, ohne Internet; schnell und mit gutem Datenschutz. & Qi2 \tabularnewline
Neuer Standard für kabelloses Laden & Kabelloses Laden mit magnetischer Ausrichtung, ähnlich MagSafe – bequemer aufladen. & eSIM \tabularnewline
Eingebettete SIM & SIM ohne physische Karte, direkt ins Smartphone geschrieben; ein Anbieterwechsel erfordert keinen Kartenwechsel. & LTPO \tabularnewline
Adaptive Bildwiederholrate-Technik & Die Bildwiederholrate passt sich automatisch an den Inhalt an (1–120 Hz) – flüssig und stromsparend. & Anhang D: Wo findet man verlässliche Ressourcen? \tabularnewline
Offizieller Support hat Vorrang: & Apple-Nutzer & : apple.com/cn oder iCloud.com. \tabularnewline
Android-Nutzer & : Auf den Websites der Smartphone-Marken gibt es einen Bereich "Service \& Support". & Allgemeine Werkzeuge \tabularnewline
: Apple nutzt iTunes, Huawei die Huawei-Telefon-Assistent-Software, Xiaomi die Xiaomi-Wechsel-App. & Praktische Websites und Blogger: & Tipps nachschlagen \tabularnewline
\hline
\end{tabularx}
\end{center}
